\subsection{Cancellation between the four real sectors}
Section~\ref{sec:comp-real-kernel-collection} has already removed the
$b_\pm$ pole family directly from the common integrand. The following
identities also recover that cancellation from the separately integrated
sectors and reduce the remaining functions.
The common primitive becomes shorter when the rational, $W$,
energy-logarithm and Bose-restoring contributions are combined before
reducing individual functions. Denote these primitives by
$F_{R_0}$, $F_W$, $F_{\log}$ and $F_D$, respectively. The factor
$\ln s$ multiplying the energy-logarithm primitive is included in
$F_{\log}$. All $G$ functions in this subsection have the same
basepoint $r=\sqrt{s}$ and the same final argument $u$.
The subscript $R_0$ denotes the rational kernel. Taking its integrated
finite part with the same cutoff $z>\zeta$ as in the sector assembly gives
\begin{equation}
 F_{R_0}^z=F_{\mathrm{rat}}^z.
\end{equation}
This is a change of notation for the same finite part, with no
additional matching shift or normalization factor.

Consider the pair of letters
\begin{equation}
 b_\pm=-1\pm\sqrt{1-s},\qquad
 \sqrt{1-s}=\ii\sqrt{s-1},\qquad s>1.
\end{equation}
Their complete contribution to the rational primitive is
\begin{align}
 \left.F_{R_0}\right|_{b_\pm}
 &=c_b(s)\sum_{b=b_+,b_-}
       \bigl[G(b,0;u)-G(b,-s;u)+G(b,-1;u)\bigr],\\
 c_b(s)&=\frac{4s\left(s^5+9s^4+116s^3+516s^2+307s-693\right)}
                   {3(s-1)(s+3)^3}.
\end{align}
The independently integrated Bose-restoring primitive has exactly the
opposite six coefficients:
\begin{equation}
 \left.F_D\right|_{b_\pm}=-\left.F_{R_0}\right|_{b_\pm},
 \qquad \left.F_W\right|_{b_\pm}=\left.F_{\log}\right|_{b_\pm}=0.
 \label{eq:comp-real-four-b-cancel}
\end{equation}
Thus neither letter is needed in the combined primitive. This is a
cancellation of identical functions with identical paths, so it introduces
no change of branches or boundary constants.

The letters $r$ and $-r$ cancel between the rational and $W$ sectors.
Define
\begin{equation}
 Q_c(u)=G(c,0;u)+G(c,1;u)-G(c,s;u).
\end{equation}
Their coefficients are
\begin{align}
 \left.F_{R_0}\right|_{\pm r}
 &=-\frac{16(r-1)^2}{r+1}Q_r(u)
       +\frac{16(r+1)^2}{r-1}Q_{-r}(u),\\
 \left.F_W\right|_{\pm r}&=-\left.F_{R_0}\right|_{\pm r},
 \qquad \left.F_{\log}\right|_{\pm r}=\left.F_D\right|_{\pm r}=0.
\end{align}
Combining all four sectors therefore removes twelve weight-two words
with these four first letters. The remaining alphabet is
\begin{equation}
 \left\{0,1,-1,s,-s,1+\sqrt{1-s},1-\sqrt{1-s},q_+,q_-\right\},
 \qquad q_\pm=\frac{s-3\pm\sqrt{(s-1)(s-9)}}2.
\end{equation}
The $q_\pm$ pair has the alternative representation derived above.
It is not necessary to expand every remaining $G$ function to make
these cancellations explicit.

There are 44, 19, 14 and 28 nonzero GPL terms in the four separate
primitive dictionaries. After collecting their coefficients, 37 distinct
words remain. Three shuffle reductions leave 34 words and three products
of weight-one functions, together with the rational term. The coefficient
collection uses $H=s-1$ and $r^2=s$ before identifying equal letters.
Both the change of variables in each coefficient and the shuffle
reconstruction are checked as exact algebraic identities.
This collection precedes the finite endpoint operation. The lower-limit
rational constants, the moving-boundary finite shifts, and the soft
contribution must still be combined to obtain a compact total cross
section.

\paragraph{The finite endpoint after collection.}
Let $F(u)=F_{R_0}(u)+F_W(u)+F_{\log}(u)+F_D(u)$ and
$b=s-u$. The collected primitive has only two rational coefficients
with a pole at $b=0$: its constant coefficient $c_{\varnothing}$ and
its coefficient $c_0$ multiplying $G(0;u)=\ln(u/\sqrt{s})$.
Their residues are
\begin{align}
 c_{\varnothing}(s-b)&=\frac{d_0(s)}b+O(1),
 &d_0(s)&=\frac{8s(s^3+6s^2-s-2)}3,\\
 c_0(s-b)&=\frac{d_1(s)}b+O(1),
 &d_1(s)&=-\frac{32s^2(s-5)}{3(s-1)}.
\end{align}
All other function coefficients are regular at $u=s$.
The product with the first-order expansion of $G(0;u)$ is essential:
\begin{equation}
 G(0;s-b)=\tfrac12\ln s-\frac bs+O(b^2).
\end{equation}
It contributes the finite term
\begin{equation}
 T_{\mathrm{end}}(s)=-\frac{d_1(s)}s
                 =\frac{32s(s-5)}{3(s-1)}.
 \label{eq:comp-real-four-endpoint-Taylor}
\end{equation}
Thus substituting $u=s$ in the regular coefficients alone would miss a
finite contribution. The coefficient of the double pole in the
integrand is the coefficient of $1/b$ in its primitive:
\begin{equation}
 D_4(s)=d_0(s)+\tfrac12d_1(s)\ln s.
 \label{eq:comp-real-four-double-pole}
\end{equation}

At the lower limit, every GPL vanishes. The coefficients of the
products of two weight-one functions have at most a simple pole at
$u=\sqrt{s}$, whereas the products vanish quadratically, so these terms
also give zero. The rational term has the finite limit
\begin{equation}
 F(\sqrt{s})=-\frac{2(s+1)(s^4-3s^3-41s^2+27s-8)}{3(s-1)}.
 \label{eq:comp-real-four-lower-limit}
\end{equation}
This combines the four lower limits before evaluating any polylogarithm.

To state the remaining change of cutoff, define $A_4(s)$ by the
integrand expansion
\begin{equation}
 F'(s-b)=\frac{D_4(s)}{b^2}+\frac{A_4(s)}b+O(\ln b).
\end{equation}
In particular the primitive contains $D_4/b-A_4\ln b$.
Writing $\mathcal E_4(s)$ for its finite coefficient, including
\eqref{eq:comp-real-four-endpoint-Taylor}, the four-sector finite part
in the original $z$ convention is
\begin{equation}
 F_4^z=\mathcal E_4(s)-F(\sqrt{s})
       +A_4(s)\ln\frac{s-1}{s}-\frac{D_4(s)}{s(s-1)}.
 \label{eq:comp-real-four-zfinite}
\end{equation}
Indeed, $z=(s-1)b/s-b^2/s^2+O(b^3)$ gives
\begin{equation}
 \frac1b=\frac{s-1}{sz}-\frac1{s(s-1)}+O(z),
 \qquad \ln b=\ln z+\ln\frac{s}{s-1}+O(z).
\end{equation}
The second-order boundary relation supplies the last finite term in
\eqref{eq:comp-real-four-zfinite}. The finite shifts from the $z$
convention to the common real-emission convention are subsequently
added as in \eqref{eq:comp-matching-short}.
\paragraph{Removing a conjugate pair at the endpoint.}
The remaining letters $a_\pm=1\pm\ii\sqrt{s-1}$ occur only in the
combination
\begin{equation}
 \mathcal A_{\mathrm{pair}}(u)=\sum_{a=a_+,a_-}
       \bigl[G(a,0;u)+G(a,1;u)-G(a,s;u)\bigr].
\end{equation}
Their common coefficient in the four-sector primitive is
\begin{equation}
 \gamma_a(s)=-\frac{32s(s^2-3s-1)}{(s-1)^2}\ln s
 -\frac{4(s^5+s^4-25s^3+47s^2-12s+12)}{3s(s-1)}.
\end{equation}
It is shorter to integrate the summed logarithmic derivative than to
evaluate the six complex-letter GPLs separately. Put
\begin{equation}
 B(u)=\frac{u(u-1)}{s-u},\qquad
 P(u)=u^2-2u+s=(u-a_+)(u-a_-).
\end{equation}
For $\sqrt{s}<u<s$, $B(u)>0$ and $B(\sqrt{s})=1$.
The second-letter combination is exactly $\ln B(u)$, so
\begin{equation}
 \mathcal A_{\mathrm{pair}}'(u)=\frac{P'(u)}{P(u)}\ln B(u).
\end{equation}
The factorization $P(u)=(s-u)(1+B(u))$ gives
\begin{equation}
 \frac{P'(u)}{P(u)}=\frac{B'(u)}{1+B(u)}+\frac1{u-s}.
\end{equation}
Since the derivative of
$\ln B\ln(1+B)+\operatorname{Li}_2(-B)$ is
$\ln B\,B'/(1+B)$, the primitive is
\begin{equation}
 \begin{aligned}
 \mathcal A_{\mathrm{pair}}(u)={}&\ln B(u)\ln(1+B(u))+\operatorname{Li}_2(-B(u))
                      +\frac{\pi^2}{12}\\
 &+G(s,0;u)+G(s,1;u)-G(s,s;u).
 \end{aligned}
\end{equation}
The constant follows from $\mathcal A_{\mathrm{pair}}(\sqrt{s})=0$ and
$\operatorname{Li}_2(-1)=-\pi^2/12$. All arguments in the explicit
dilogarithm term are negative real, fixing its branch throughout the
integration interval.

The separate terms contain logarithms at the upper endpoint, but their
sum has a finite limit. With $r=\sqrt{s}$ and $b=s-u$, the required
expansions are
\begin{align}
 \ln B\ln(1+B)+\operatorname{Li}_2(-B)
 &=\tfrac12\ln^2\!\left(\frac{s(s-1)}b\right)
      -\frac{\pi^2}{6}+O(b|\ln b|),\\
 G(s,0;s-b)&=\ln r\ln\frac b{s-r}
           +\operatorname{Li}_2(1-1/r)+O(b),\\
 G(s,1;s-b)&=\ln(r+1)\ln\frac b{s-r}
           +\operatorname{Li}_2\!\left(\frac r{r+1}\right)+O(b),\\
 G(s,s;s-b)&=\tfrac12\ln^2\frac b{s-r}.
\end{align}
Both powers of $\ln b$ cancel exactly. The limit first takes the form
\begin{equation}
 \mathcal A_{\mathrm{pair}}(s)=\operatorname{Li}_2(1-1/r)
 +\operatorname{Li}_2\!\left(\frac r{r+1}\right)
 +\tfrac12\ln^2[r(r+1)]-\frac{\pi^2}{12}.
\end{equation}
The reflection and inversion identities of the dilogarithm then give
\begin{equation}
 \boxed{\mathcal A_{\mathrm{pair}}(s)=\frac{\pi^2}{4}
       +\operatorname{Li}_2(-1/\sqrt{s})
       -\operatorname{Li}_2(1/\sqrt{s})
       +\frac12\ln s\ln(s-1).}
 \label{eq:comp-real-a-pair-endpoint}
\end{equation}
For example, use
$\operatorname{Li}_2(r/(r+1))=-\operatorname{Li}_2(-r)
-\tfrac12\ln^2(r+1)$ and
$\operatorname{Li}_2(-r)+\operatorname{Li}_2(-1/r)
=-\pi^2/6-\tfrac12\ln^2r$, together with the reflection identity
for $\operatorname{Li}_2(1-1/r)$.
The final two dilogarithm arguments lie in $(-1,1)$ for every $s>1$.
Thus the complete contribution is
$\gamma_a(s)\mathcal A_{\mathrm{pair}}(s)$, and the six complex-letter endpoint values
are unnecessary.
